% interaction/foundations.tex — Spec, Transcript, Decoration, Strategy

\section{Foundations: Dependently-Typed Interaction}\label{sec:interaction-foundations}

The interaction core of ArkLib is designed to model protocols in which the
remainder of the interaction may depend on the transcript so far.  This extra
generality is not introduced for its own sake.  It is needed to represent
protocols where an early message determines the later message space, ambient
field, or even the number of remaining rounds, while still supporting
transcripts, composition, oracle semantics, and security reasoning in one
uniform framework.  We begin with the modeling problem, then define the core
interaction tree, and finally explain how additional metadata is layered over
that tree rather than baked into a bespoke inductive.

\subsection{Motivation: from flat indexing to dependent trees}

Any general framework for interactive proofs must allow later parts of the
protocol to depend on what has already been said.  In many familiar examples the
round structure is fixed in advance, so this dependence is easy to ignore.  But
there are important protocols in which the transcript changes the ambient object
of later interaction: the message space, the field, or the remaining length of
the protocol itself.

A natural first attempt at modeling an $n$-round protocol is therefore to fix
the number of rounds in advance and store the move types in a vector indexed by
$\Fin\;n$:
%
\[
  \ProtocolSpec\;n \;\defeq\;
    \bigl\{
      \mathit{dir} : \Fin\;n \to \Direction,\;\;
      \mathit{type} : \Fin\;n \to \Type
    \bigr\}.
\]
%
This is the approach taken in the original ArkLib core.  It works for protocols
whose round types are fixed statically, but it presupposes that the schedule of
interaction is known before the protocol begins.  That is already too rigid for
protocols where a later challenge space depends on an earlier message, or where
an early claim determines how many rounds remain.

Even in the fixed-round setting, the global index creates a second problem:
composition and prefixes become dominated by index arithmetic.  Accessing the
transcript up to round~$k$ requires casting from $\Fin\;k$ to $\Fin\;n$;
composing two protocols of lengths $m$ and $n$ forces every index
$i : \Fin\;(m+n)$ to be split into cases $i < m$ and $i \ge m$, each requiring
dependent type casts with ad-hoc arithmetic lemmas.  In practice, composition
definitions accumulate several layers of \texttt{Fin.castLE},
\texttt{Fin.castSucc}, and related casts that make proofs brittle and
definitions hard to read.  This implementation pain is a symptom of the deeper
semantic mismatch.

An intermediate improvement replaces the $\Fin\;n$-indexed vector with an
inductive list of $(\Direction \times \Type)$ pairs.  This removes the
cast burden, because structural recursion on a list naturally decomposes
``first round'' from ``remaining rounds.''  All $\Fin$-arithmetic disappears.
But the list model still fixes the entire schedule in advance: it improves the
recursion principle without changing what can be expressed.

However, the list-based model still fixes every round's type \emph{independently
of prior moves}.  For most standard protocols---sumcheck, FRI, and their
variants---this is adequate, because message types are either constant or
parameterized by data external to the protocol (e.g., a fixed field~$\F$).
But there are protocols where genuine inter-round type dependence is
mathematically forced:

\begin{itemize}
  \item \textbf{Random-modulus protocols (Zinc~\cite{Zinc}).}
    The prover and verifier begin with a relation over~$\Q$.
    The verifier samples a random prime~$q$; all subsequent rounds
    operate over~$\F_q$.  The type of later messages---field elements
    modulo~$q$---depends on the verifier's choice.  The transcript does not
    merely affect the values exchanged later; it changes the message space
    itself.

  \item \textbf{Variable-length protocols (zkVMs such as Jolt~\cite{Jolt}).}
    The prover commits a claimed execution trace of length~$T$.
    Subsequent sumcheck invocations require $\log_2 T$ rounds, so the
    \emph{number of rounds} depends on a prior message.  More generally,
    parameters such as RAM size, polynomial layout, or commitment structure
    may be sent as early messages and shape the rest of the protocol.
\end{itemize}

Because ArkLib aims to model \emph{all} IOP-based protocols, including these,
we adopt a \emph{dependently-typed specification of interaction}, where each
round's continuation may depend on the move actually played.  The right object
is therefore not a list of round slots, but a well-founded tree of possible next
moves.

\subsection{Interaction specifications and transcripts}

\begin{definition}[Interaction Specification]
  \label{int:spec}
  An \emph{interaction specification} is a well-founded tree defined by the
  inductive type
  \[
    \mathsf{Spec} \;\defeq\;
    \begin{cases}
      \mathsf{done} \\
      \mathsf{node}\;(X : \Type)\;(\mathit{rest} : X \to \mathsf{Spec})
    \end{cases}
  \]
  Each internal node carries a type~$X$ of moves that can be exchanged;
  the remaining protocol $\mathit{rest}\;x$ may depend on the chosen
  move~$x$.  Terminal nodes $\mathsf{done}$ indicate the end of the
  interaction.
  \lean{Interaction.Spec}
\end{definition}

One should read $\mathsf{Spec}$ as a protocol tree.  Each internal node records
the next \emph{move space}~$X$, and each move $x : X$ determines the remainder
of the protocol.  Fixed-round protocols are recovered as a special case; the
non-dependent list encoding appears via
$\mathsf{Spec.ofList} : \List\;\Type \to \mathsf{Spec}$.

\begin{definition}[Transcript]
  \label{int:transcript}
  A \emph{transcript} of a $\mathsf{Spec}$ is a complete root-to-leaf path
  through the tree: at each node, a concrete move is recorded.
  \[
    \Transcript : \mathsf{Spec} \to \Type, \qquad
    \Transcript\;\mathsf{done} = \Unit, \qquad
    \Transcript\;(\mathsf{node}\;X\;\mathit{rest}) =
      (x : X) \times \Transcript\;(\mathit{rest}\;x).
  \]
  \lean{Interaction.Spec.Transcript}
\end{definition}

This is the point at which the formal object should feel natural to both
audiences: cryptographically, a transcript is simply the record of all messages
actually exchanged; type-theoretically, it is a dependent path through the
interaction tree.

For readers from type theory, $\mathsf{Spec}$ is a $W$-type in the
Hancock--Setzer sense~\cite{HancockSetzer2000}.  It can also be viewed as a
well-founded dependent game tree, closely related to the history-dependent game
trees studied by Escard\'o--Oliva~\cite{EscardoOliva2023}.  We mention these
connections because they explain why the subsequent structure composes so
smoothly, not because they are prerequisites for reading the rest of this
chapter.

\subsection{Decorations as displayed algebras}

On its own, a $\mathsf{Spec}$ says nothing about \emph{who} makes each move,
\emph{how} moves are computed, or what oracle interfaces are available.  Rather
than baking this metadata into a separate inductive type for each concern---which
would force us to duplicate all transcript and composition infrastructure---we
layer it as a \emph{decoration}.

\begin{definition}[Decoration]
  \label{int:decoration}
  Given a type family $S : \Type \to \Type$, a \emph{decoration}
  $\mathsf{Decoration}\;S\;\mathit{spec}$ attaches an $S\;X$ label at each
  internal node with move type~$X$:
  \[
    \mathsf{Decoration}\;S\;\mathsf{done} = \Unit, \qquad
    \mathsf{Decoration}\;S\;(\mathsf{node}\;X\;\mathit{rest}) =
      S\;X \times \textstyle\prod_{x : X}\,
        \mathsf{Decoration}\;S\;(\mathit{rest}\;x).
  \]
  \lean{Interaction.Spec.Decoration}
\end{definition}

Decorations admit a natural transformation $\mathsf{map} :
(\forall\,X,\; S\;X \to T\;X) \to \mathsf{Decoration}\;S \to
\mathsf{Decoration}\;T$ satisfying the expected functoriality laws.

\begin{definition}[Displayed Decoration (Refine)]
  \label{int:decoration-refine}
  A \emph{displayed decoration} or \emph{refinement}
  $\mathsf{Decoration.Refine}\;F\;\mathit{spec}\;d$ is a decoration fibered
  over an existing decoration~$d$: at each node with label $l : L\;X$ from~$d$,
  it attaches data in $F\;X\;l$.
  \lean{Interaction.Spec.Decoration.Refine}
  \uses{int:decoration}
\end{definition}

The key benefit is practical.  Because roles, oracle interfaces, and other
metadata are all instances of $\mathsf{Decoration}$ (or
$\mathsf{Decoration.Refine}$), every operation on $\mathsf{Spec}$---transcripts,
append, replicate, state chains---is defined \emph{once} and then reused at
every metadata layer.

For readers from PL and type theory, this is the \emph{displayed algebra}
pattern emphasized by McBride~\cite{McBride2010}, together with the
\emph{ornament} perspective developed further by
Dagand--McBride~\cite{DagandMcBride2014}.  In this chapter, however, the main
point is simpler: decorations let us enrich the same interaction tree without
rebuilding the entire framework for each new concern.

\subsection{Strategies}

\begin{definition}[Strategy]
  \label{int:strategy}
  A \emph{strategy} $\Strategy\;m\;\mathit{spec}\;\mathit{Output}$ plays
  through a $\mathsf{Spec}$, choosing moves and interleaving monadic effects
  in~$m$, producing a transcript-dependent result:
  \[
    \Strategy\;m\;\mathsf{done}\;\mathit{Output}
      = \mathit{Output}\;\langle\rangle, \qquad
    \Strategy\;m\;(\mathsf{node}\;X\;\mathit{rest})\;\mathit{Output}
      = (x : X) \times m\bigl(\Strategy\;m\;(\mathit{rest}\;x)\;
        (\lambda\,p.\;\mathit{Output}\;\langle x,p\rangle)\bigr).
  \]
  \lean{Interaction.Spec.Strategy}
\end{definition}

The definition is by structural recursion on $\mathsf{Spec}$, following the
Hancock--Setzer pattern.  This avoids the positivity issues that arise when
defining strategies as a coinductive free monad over a generic~$m$.

Cryptographically, a strategy is the notion of protocol algorithm used
throughout the framework: honest provers, honest verifiers, oracle simulators,
and related objects are all instances obtained by choosing an appropriate monad
and output family.

$\mathsf{Strategy.run}$ executes a strategy, returning the full transcript
and the dependent output.  $\mathsf{Strategy.mapOutput}$ is the functorial
action on the output family, with $\mathsf{mapOutput\_id}$ and
$\mathsf{mapOutput\_comp}$ establishing that it forms a functor.
